\section{Conclusion}\label{sec:conclusion}

We measured how much of a retrieval gain reaches the answer when a RAG system swaps one open embedding model for another. On single-hop NQ, between \ev{T.lnq.R4} and \ev{T.lnq.R3} EM points followed each Hit@5 point gained over MiniLM-L6, for both readers (RQ1). On multi-hop HotpotQA the ratio was larger, mainly because better encoders more often retrieved the second supporting passage, which Hit@5 does not count. Among dense encoders, bge-large gave a small EM gain over e5-base-v2 on HotpotQA and none detectable on NQ, and Qwen3-Embedding-0.6B gave lower EM (RQ2). Top-ranked non-gold passages lowered EM, gold position did not change it detectably, a reranker helped only the weak BM25 retrieval stage, and ten passages narrowed the gaps between retrievers (RQ3). Encoding throughput differed \ev{E7.ratio.R1R6}-fold, while approximate search and the sub-corpus did not change the order of retrievers (RQ4).

Within this setup, retrieval benchmarks overstate the answer value of stronger encoders on single-hop questions, and Hit@k understates it on multi-hop questions. We recommend that RAG studies report answer gains next to retrieval gains, with paired intervals, and use retrieval metrics that require all evidence a question needs. Future work should test larger readers, other domains and a full corpus for the readers, and should replace EM with validated human or model judgements.
